\documentclass[12pt]{article}
\usepackage[utf8]{inputenc}
\usepackage[spanish]{babel}
\usepackage{amsmath}
\usepackage{amsfonts}
\usepackage{amssymb}
\usepackage{geometry}
\geometry{letterpaper, margin=2.5cm}

\begin{document}

\begin{center}
    \Large \textbf{Evaluación de Examen 1: Unidad I} \\
    \large Física para Ciencias de la Salud (005-1713) \\
    \vspace{0.5cm}
    \textbf{Estudiante:} Danebsas Ramos \quad \textbf{C.I.:} 34.254.529
    \vspace{0.3cm} \\
    \large \textbf{Calificación Final: 8/10 puntos}
\end{center}

\vspace{0.5cm}

\section*{Análisis de Resultados}

\subsection*{1. Ángulo plano (Conversión a radianes)}
\textbf{Procedimiento y resultado:} Aplica correctamente el factor de conversión formal $\left(\frac{\pi \text{ rad}}{180^\circ}\right)$. Simplifica la fracción adecuadamente para obtener la expresión exacta e irreductible de $\frac{5\pi}{12}$ radianes. Un procedimiento limpio y directo. \\
\textbf{Veredicto:} \textbf{Correcto} (2/2 pts).

\subsection*{2. Sistemas de coordenadas (Longitud del hueso)}
\textbf{Procedimiento y resultado:} Excelente planteamiento y resolución. La estudiante construye un gráfico sumamente claro donde identifica el triángulo rectángulo en el plano. Aplica el Teorema de Pitágoras $d^2 = (6\text{ cm})^2 + (8\text{ cm})^2$, resolviendo paso a paso hasta llegar al resultado correcto de $10\text{ cm}$. \\
\textbf{Veredicto:} \textbf{Correcto} (2/2 pts).

\subsection*{3. Vectores (Fuerza resultante)}
\textbf{Procedimiento y resultado:} Presenta un error conceptual respecto al enunciado. El problema solicitaba encontrar la fuerza resultante, lo cual corresponde a la suma vectorial ($\vec{F}_R = \vec{F}_1 + \vec{F}_2$). La estudiante, sin embargo, plantea y ejecuta una diferencia de vectores ($\Delta\vec{F} = \vec{F}_2 - \vec{F}_1$). Aunque la operación matemática de la resta fue ejecutada correctamente obteniendo $-60\hat{i} + 10\hat{j}$ N, no responde a lo que se pedía en el problema (que era $30\hat{i} + 60\hat{j}$ N). \\
\textbf{Veredicto:} \textbf{Parcialmente correcto} (Manejo adecuado del álgebra vectorial, pero error en el planteamiento físico) (1/2 pts).

\subsection*{4. Potencias de diez (Notación científica y prefijo)}
\textbf{Procedimiento y resultado:} ¡Excelente! La estudiante no solo expresa correctamente el número en notación científica ($1,25 \times 10^{-5}$ m), sino que además completa satisfactoriamente la segunda parte de la pregunta indicando el nombre del prefijo del Sistema Internacional equivalente al desplazar la coma ("$12,5$ micro metro"). \\
\textbf{Veredicto:} \textbf{Correcto} (2/2 pts).

\subsection*{5. Sistemas de unidades (Conversión de densidad)}
\textbf{Procedimiento y resultado:} Plantea los factores de conversión en cadena de manera magistral, indicando claramente las equivalencias de masa ($1 \text{ kg} = 1000 \text{ g}$) y volumen ($1 \text{ m}^3 = 1.000.000 \text{ cm}^3$). Lamentablemente, el procedimiento se detiene al plantear la ecuación final y no se efectúa el cálculo aritmético que daría el resultado final de $1030 \text{ kg/m}^3$. \\
\textbf{Veredicto:} \textbf{Incompleto / Parcialmente correcto} (Planteamiento excelente, falta de resultado final) (1/2 pts).

\vspace{0.5cm}
\section*{Comentarios Generales y Sugerencias}
El examen de la estudiante refleja un gran nivel de organización, gráficos impecables y un muy buen razonamiento. Fue la única capaz de resolver íntegramente la pregunta de los prefijos del S.I. 
\begin{itemize}
    \item Se sugiere repasar el concepto de "fuerza resultante" en estática y dinámica, recordando que corresponde a la sumatoria de todas las fuerzas actuantes.
    \item Recordar siempre completar los cálculos algebraicos hasta expresar el resultado final de la variable que se está buscando (como ocurrió en la Pregunta 5).
\end{itemize}

\end{document}